\chapter{Grundlagen}

\section{Grundlagen der Audio-DSP}

\subsection{Physikalische Grundlage des Audiosignals}

In der Audiosignalverarbeitung wird die Erfassung, Speicherung, Anzeige und Erzeugung von Audiosignale gesprochen. Ebenso als Extraktion und Neu-Codierung von Informationen.
In der Regel entsteht durch mechanische Schwingungen in der Luft ein Schall mit einer Schallgeschwindigkeit $c \approx 343 \, \text{m/s}$ im Raum ausbreitet. Um den Schalldruck als Signal zu erfassen ist die einfachste Darstellung die Sinuswelle. \cite{rossing2007springer}

\begin{equation}
	x(t) = A \cdot \sin(\omega t + \phi)
\end{equation}

Die wichtigsten Größen im Zeitbereich sind:

\begin{itemize}
	\item \textbf{Amplitude} $A$: die Größe der Schwingung, also die maximale Druckabweichung vom Ruhezustand.
	\item \textbf{Frequenz} $f$: die Anzahl der Schwingungswiederholungen pro Sekunde, gemessen in Hertz~(Hz).
	\item \textbf{Periode} $T$: die Dauer einer vollständigen Schwingung, mit
	\begin{equation}
		T = \frac{1}{f}
	\end{equation}
	\item \textbf{Phase} $\phi$: die Anfangslage der Schwingung im Zyklus, angegeben in Radiant.
\end{itemize}

\subsection{Abtastung und Quantisierung}
Das Mikrofon wandelt den Schalldruck $p(t)$ in eine analoge elektrische Spannung $u(t)$ um - ein kontinuierliches Signal, das anschließend durch den ADC abgetastet und quantisiert wird. Damit ist die Verbindung zur digitalen Signalverarbeitung hergestellt. Dabei ist Aliasing ein zentrales Problem,   \cite{rossing2007springer}.


\subsection{Digitale Audiosignale als Sample-Folgen}

\subsection{Signal-Rausch-Verhältnis (SNR)}
Wichtig für den späteren Verständnis ist das Signal-Rausch-Verhältnis. SNR beschreibt das Verhältnis der Nutzleistung eines Signals zur Leistung des überlagerten Rauschens. Je höher das SNR, desto weniger wird das Nutzsignal durch Rauschen beeinträchtigt.

Das SNR ist definiert als:

\begin{equation}
	\text{SNR} = \frac{P_{\text{Signal}}}{P_{\text{Rauschen}}}
\end{equation}

Da Leistungsverhältnisse sehr große oder sehr kleine Zahlen annehmen können, 
wird das SNR in der Praxis logarithmisch in Dezibel (dB) angegeben:

\begin{equation}
	\text{SNR}_{\text{dB}} = 10 \cdot \log_{10} 
	\left( \frac{P_{\text{Signal}}}{P_{\text{Rauschen}}} \right)
\end{equation}

Sind statt Leistungen die Amplituden bekannt, gilt:

\begin{equation}
	\text{SNR}_{\text{dB}} = 20 \cdot \log_{10} 
	\left( \frac{A_{\text{Signal}}}{A_{\text{Rauschen}}} \right)
\end{equation}

Typische Orientierungswerte für die Sprachverarbeitung:

\begin{itemize}
	\item $\text{SNR} < 0\,\text{dB}$: Rauschen lauter als Nutzsignal
	\item $\text{SNR} \approx 10\,\text{dB}$: Sprache gerade noch verständlich
	\item $\text{SNR} \approx 20\,\text{dB}$: gute Sprachqualität
	\item $\text{SNR} > 40\,\text{dB}$: Studioqualität
\end{itemize}

Für die spektrale Subtraktion ist das SNR die zentrale Bewertungsgröße: 
Ziel des Algorithmus ist es, das SNR des verrauschten Eingangssignals 
$y[n] = x[n] + d[n]$ durch Unterdrückung des Rauschspektrums $D[k]$ 
zu maximieren, ohne dabei das Nutzsignal $x[n]$ wesentlich zu verzerren. \cite{Loizou2013}

\section{Frequenzanalyse mit DFT und FFT}

\subsection{Zeitbereich und Frequenzbereich}

\subsection{Diskrete Fourier-Transformation}

\subsection{Fast Fourier Transformation}

\subsection{Amplituden- und Leistungsspektrum}

\subsection{Frequenzauflösung}

\section{Kurzzeitanalyse und STFT}

\subsection{Nichtstationarität von Sprache}

\subsection{Frame-basierte Verarbeitung}

\subsection{Short-Time Fourier Transform}

\subsection{Magnitude und Phase}

\section{Fensterfunktionen}

\subsection{Blockbildung und Randdiskontinuitäten}

\subsection{Spectral Leakage}

\subsection{Hanning- und Hamming-Fenster}

\subsection{Bedeutung für spektrale Subtraktion}

\section{Overlap-Add und Rekonstruktion}
\subsection{Inverse FFT}
\subsection{Überlappende Frames}
\subsection{Rekonstruktion des Zeitsignals}
\section{Echtzeitanforderungen}
\subsection{Latenz}
\subsection{Puffergröße}
\subsection{Rechenzeit pro Frame}
\subsection{Overrun und Underrun}
\subsection{Echtzeitfähigkeit der spektralen Subtraktion}
\section{Zusammenfassung}
